\SourceLeaf{337}%
\(T^0\) ansetzen) und rechnet nach den vier Formeln
\[
\begin{aligned}
\beta&=(2)r\cos T\\
B&=B^0-\tfrac12\beta\\
\tau&=(1)r\sin T\operatorname{tang}B\\
T&=T^0-\tfrac12\tau
\end{aligned}
\]
Den Wert von (2) sollte man der Tafel mit dem Argument \(B\) entnehmen. Beim ersten Mal kann man das Argument \(B^0\) verwenden, sofern man nicht glaubt, einen genaueren Wert von \(B\) durch Schätzung vorwegnehmen zu können. Den Wert von (1) entnimmt man der Tafel mit dem gerade gefundenen Wert von \(B\).

Man wiederholt dieselbe Rechnung mit dem durch die vierte Gleichung gefundenen Wert von \(T\) und entnimmt (1) und (2) mit dem bereits verbesserten \(B\) der Tafel. Falls nötig, wiederholt man die Rechnung erneut, bis das Ergebnis stabil bleibt, also bis man durch die vierte Formel denselben Wert von \(T\) erhält, von dem man zuletzt ausgegangen ist. Für alle diese Rechnungen verwendet man nur fünfstellige Logarithmen.

Bei den weiteren Wiederholungen rechnet man mit siebenstelligen Logarithmen, bezieht die logarithmischen Korrekturen von \(\log\tau\) und \(\log\beta\) mit ein und setzt \(B=B^0-\tfrac12b,\ T=T^0-\tfrac12t\). Erst wenn auch diese Rechnung stabile Ergebnisse geliefert hat, berechnet man \(\lambda\) und \(l\) nach den Formeln am Schluss von Artikel 33. Zur Erläuterung dieser Vorschriften folgen hier die wichtigsten Schritte eines Beispiels, das wie oben in Artikel 20 bei der Rechnung auf der Kugel von der Dreiecksseite Brocken-Inselsberg ausgeht.

Bei der Rechnung auf dem Ellipsoid beträgt die Breite des Brockens
\[
=51^0\ 48'\ 1''\ 9294=B^0
\]
und das Azimut der Seite Brocken-Inselsberg beträgt
\[
=5^0\ 42'\ 21''\ 7699=T^0
\]
Der Logarithmus der Dreiecksseite in Toisen stimmt bis zur siebten Dezimalstelle mit dem in der konformen Abbildung auf der Kugeloberfläche überein, nämlich \(=4{,}7353929\). In der Einheit, die unserer Hilfstafel zugrunde liegt, ist somit \(\log r=5{,}0251757\).
